\chapter{Eje sintético}\label{anx:sintetico}

Elliptic no publica el significado de sus variables ni anota qué patrón de lavado sigue cada
transacción ilícita, de modo que sobre él solo se puede medir la estabilidad de las
explicaciones, no si señalan lo correcto. Para observar esa segunda propiedad se construyó un
grafo sintético de transacciones en el que cada transacción ilícita pertenece a una tipología
conocida y en el que se sabe qué aristas y qué variables la delatan. Este anexo resume sus
resultados como evidencia complementaria; no forma parte del contraste de hipótesis.

\section{Diseño}

El generador (\texttt{phase1/synthetic\_aml\_generator.py}, opción \texttt{--v4}, semilla 42)
produce un grafo con tipologías de lavado insertadas sobre un fondo de transacciones
legítimas. Se verificó que ninguna variable aislada permite separar las clases (AUC de una
sola variable entre 0,50 y 0,62) y que las cuatro arquitecturas aprenden la tipología con un
margen de al menos 0,10. Se entrenaron 144 configuraciones: las 4 arquitecturas, 4 escenarios,
3 técnicas de compensación y 3 semillas, explicadas con GNNExplainer, PGExplainer y Shapley
por muestreo.

Los escenarios del eje sintético difieren de los de Elliptic: los niveles 1:50 y 1:100 se
obtienen quitando transacciones ilícitas. Además se miden la \textbf{plausibilidad}
(coincidencia de las variables o aristas señaladas con las que de verdad definen la
tipología; para las variables, la fracción de las variables verdaderas que aparecen entre las 3
principales) y la \textbf{fidelidad} (cuánto cambia la predicción al retirar lo señalado, fid+,
o al conservar solo lo señalado, fid$-$), cada una frente a una explicación al azar del mismo
tamaño. La fidelidad se calcula sobre lo que cada explicador explica: retirando aristas en
GNNExplainer y PGExplainer, y retirando variables en Shapley. Por eso la fidelidad de
GNNExplainer y la de Shapley no son comparables entre sí; cada una solo se compara con su propio
valor al azar.

\section{Resultados}

\begin{table}[ht]
\centering
\caption{Eje sintético: estabilidad (Spearman global), plausibilidad de variables y fidelidad
(fid+; sobre aristas en GNNExplainer y sobre variables en Shapley), con su valor al azar entre
paréntesis. Medias sobre 36 configuraciones por fila.
Fuente: \texttt{results\_phase1\_v4/summary\_v4.csv}.}
\label{tab:sintetico}
\small
\begin{tabular}{llccc}
\toprule
Arquitectura & Explicador & Estabilidad & Plausibilidad de variables & Fidelidad (fid+) \\
\midrule
GCN & GNNExplainer & 0,848 & 0,296 (0,150) & 0,147 (0,030) \\
    & Shapley      & 0,984 & 0,114 (0,150) & 0,195 (0,063) \\
GraphSAGE & GNNExplainer & 0,907 & 0,404 (0,150) & 0,363 (0,048) \\
          & Shapley      & 0,953 & 0,157 (0,150) & 0,135 (0,092) \\
GAT & GNNExplainer & 0,850 & 0,244 (0,150) & 0,264 (0,054) \\
    & Shapley      & 0,959 & 0,118 (0,150) & 0,118 (0,101) \\
TAGCN & GNNExplainer & 0,957 & 0,447 (0,150) & 0,277 (0,053) \\
      & Shapley      & 0,945 & 0,159 (0,150) & 0,154 (0,083) \\
\bottomrule
\end{tabular}
\end{table}

La Tabla \ref{tab:sintetico} muestra el resultado central del eje: el explicador más estable
no es el más correcto. Shapley por muestreo es el más estable en tres de las cuatro
arquitecturas (Spearman entre 0,95 y 0,98), pero su plausibilidad de variables no supera al
azar (entre 0,11 y 0,16 frente a 0,15). GNNExplainer es menos estable, pero es el único cuya
plausibilidad de variables supera claramente al azar (entre 0,24 y 0,45). Ambos explicadores
superan a su propia línea base de fidelidad, salvo Shapley en GAT (0,118 frente a 0,101,
$p = 0{,}22$), pero esas cifras no
permiten decir cuál es más fiel, porque miden cosas distintas. La estabilidad de GNNExplainer baja con el desbalance, de 0,92 en los niveles
nativo y 1:1 a 0,86 en 1:100, mientras que la de Shapley no cambia.

Este resultado apoya una de las conclusiones del estudio principal: la estabilidad es
necesaria pero no suficiente, y debe complementarse con pruebas de que la explicación depende
del modelo.

\section{Advertencias}

\begin{itemize}
  \item La plausibilidad de aristas no es comparable entre arquitecturas. TAGCN ve 6 saltos y
        su subgrafo receptivo contiene casi todo el grafo (mediana de 9.311 nodos y 47.311
        aristas), mientras que las otras tres ven 2 saltos (unos 36 nodos y 78 aristas). La
        plausibilidad al azar vale 0,001 en TAGCN y cerca de 0,47 en las demás, y la ganancia
        aparente de TAGCN se debe a esa línea base. Por eso la tabla no la incluye.
  \item En las arquitecturas de dos saltos, la plausibilidad de aristas de los explicadores
        apenas supera al azar (entre 0,01 y 0,03), y la de PGExplainer queda por debajo.
  \item PGExplainer falló en 24 de 432 combinaciones, todas en los escenarios 1:50 y 1:100.
\end{itemize}
